\section{Living the Dream}

\begin{quotex}
As the oyster makes its own shell, so the mind of man creates and necessitates his own life and fate.
\flright{\textsc{Ananda Coomaraswamy}, \emph{Myths of the Hindus and Buddhists}}
\end{quotex}


In Vedantic teaching, there are three normal states of consciousness: deep sleep, dreaming sleep, and waking consciousness. In addition, there is a fourth state, \emph{turiya}, which is not really a separate state, but encompasses all states. In Man and his Becoming, \textbf{Rene Guenon} refers to the Mandukya Upanishad to indicate how this state is known (note that this is really a \emph{via negativa}):

\begin{itemize}


\item It knows neither internal nor external objects in an analytical sense
\item It is not a synthetic whole of integral knowledge: it is neither knowing nor not knowing 
\item It is invisible: not perceptible by any sense 
\item It is actionless: changeless identity 
\item Incomprehensible: it comprehends all 
\item Unthinkable: it cannot be clothed in any form 
\item Indescribable: it cannot be qualified by any particular attribute 
\item It is the unique essence of the Self: Atman is present in all states 
\end{itemize}


Atman is one and simple, hence it cannot be comprised of multiple states, although the states may be appearances of Atman. Since Atman is present in all states, it cannot be any particular object of consciousness.

In the other states, what, then, is not permanent? In a dream, none of the persons, places, or events that appear in it is permanent. When I wake up, they disappear. Like the oyster making his own shell, I create all of that even if I have no awareness of having done so. The embarrassment I felt when I was caught naked out in the street was my creation. So was the anxiety experienced when stuck on a high building which was spontaneously disintegrating. Ditto for the terror caused when I was chased by a raging lion, but then fell to the ground and was unable to get up and run.

All that negativity is my own creation in the dream state. Nevertheless, it is gone the moment I awaken. Usually that moment is not under my conscious will as other factors determine when I will wake up. Occasionally, there are exceptions as I recently wrote about one. When I realized my book disappeared, I knew I was dreaming and woke myself up.

The conditions are analogous in the waking state. All the persons, places, and events, that so preoccupied me and held my attention, suddenly disappear, at least for me, at the moment of death. At that point, I finally will realize that my so-called life was not real at all, although what I wake up to may be better or worse than the waking life it replaced. The analogous question arises. If there are clues in the dreaming state that indicate that I am dreaming, and thus act as triggers to awaken me, are there also triggers in my waking state of consciousness to that serve a similar purpose? It must be possible since that is the teaching about the \emph{jivanmukti}, i.e., one who awakens while still alive in the body.

Like the oyster, I create my own life and necessitate my own fate. Do you doubt that? If so, there are two possible reasons:

\begin{enumerate}


\item you are unaware of your own true will and how it creates your life

\item you really don't know what your life is
\end{enumerate}


Your life is what remains after you die and that life is eternal. That, then, is the life you are creating and there can be no doubt about it. So what are the triggers? How many do you need to remind you to wake up? \textbf{St Augustine} says, which we quote on the front page, that the truth resides in the interior of man. \textbf{Guido De Giorgio} says that restoration must begin in the interior. Therefore, changing the persons, places, conditions, and events of our outer life will amount to nothing if we don't first address the disorder in our inner life.

It is our soul that is eternal, so it is our soul that requires our attention. How many times have we read that the true holy war is within? Do we believe it and act as if it were true, or are we more interested in our petty battles we fight in the outer world? How many times have we read that the intellect must dominate the lower forces in the soul? Do we work to develop that power of concentration and control, or are we instead obsessed with dominating people and conditions in our daily lives?

We know what to do. Undergo an intellectual realization. Quiet the randomness of the mind. Be truth tellers and Spirit lovers. Allow the Logos to dominate one's consciousness.


\flrightit{Posted on 2013-02-11 by Cologero}

\begin{center}* * *\end{center}

\begin{footnotesize}
\begin{sffamily}
\texttt{spaceloom on 2013-02-12 at 09:34 said:}

``Allow the Logos to dominate one's consciousness.'' true, but how?

\hfill

\texttt{Cologero on 2013-02-12 at 11:17 said:}

I think we've offered some suggestions for that several times, most recently in \textit{Mind Fasting}\footnote{\url{http://www.gornahoor.net/?p=5394}}.

\hfill

\texttt{Mleccha on 2013-06-11 at 23:18 said:}

I have some question:

If Turiya is the Eternal Atman, or more technically, never enter in the stream of time and space even thouth is ever-present this mean that awakening of Turyia as our existencial condition(not as a peak experience to `feel` but as any true initiatic knowlegde, our very condition of being) and non identification with any Sheaths, ie, like ramana maharshi says Self-realization being only the realization(and by definition, non identification with any gross, subtle or causal state) of ow`s true nature in guenonian terminology means Liberation/Deliverance, the Union with Brahman?

\hfill

\texttt{August on 2013-06-12 at 07:23 said:}

Yes, pretty much.


\hfill

\end{sffamily}
\end{footnotesize}

